\chapter{Filtro attivo Sallen-Key}
\section{Introduzione}

Il \textbf{filtro Sallen‑Key} è una delle topologie più comuni per la realizzazione di filtri attivi del secondo ordine.
Questi filtri utilizzano un OpAmp insieme a resistori e condensatori per ottenere
risposte in frequenza desiderate (passa‑basso, passa‑alto, passa‑banda, ecc.).

Le principali caratteristiche del filtro Sallen‑Key sono:

\begin{itemize}
    \item \textbf{Semplicità di progettazione:} è facile calcolare i valori dei componenti per ottenere una frequenza di taglio e un fattore di merito specifici.  

    \item \textbf{Alta impedenza di ingresso e bassa impedenza di uscita:} grazie all’OpAmp, il filtro non carica lo stadio precedente e può pilotare comodamente lo stadio successivo.  
          Questo rende la topologia particolarmente adatta in catene audio e in strumentazione, dove
          è fondamentale isolare i vari blocchi per evitare interazioni indesiderate.
    \item \textbf{Risposta precisa:} consente di realizzare facilmente filtri con caratteristiche standard come Butterworth, Chebyshev o Bessel.  
          La possibilità di utilizzare componenti standard (ad esempio coppie \(R_1 = R_2\) e \(C_1 = C_2\)) semplifica ulteriormente il progetto e la taratura in produzione.
\end{itemize}

Nei paragrafi seguenti verranno analizzate le implementazioni \emph{passa‑basso LP},
\emph{passa‑alto HP} e \emph{passa‑banda BP} della topologia Sallen‑Key,
con particolare attenzione alla \emph{risposta in frequenza} e alla \emph{fase}, utilizzando grafici di Bode per diversi valori di \(Q\).


\newpage
\section{Filtro Sallen-Key a guadagno unitario}


L'obiettivo dell'analisi è determinare la funzione di trasferimento $H(s) = \frac{V_{out}(s)}{V_{in}(s)}$. Per fare ciò, applichiamo la \textbf{Legge di Kirchhoff delle Correnti (LKC)} al nodo principale $v_x$. 

\SallenKeyZ
Assumendo che l'amplificatore operazionale sia ideale e operi in zona lineare, possiamo scrivere l'equazione di equilibrio delle correnti al nodo $v_x$ come segue:

\begin{equation}
    i_{Z_1} = i_{Z_2} + i_{Z_3}
    \label{eq: LKT V_x}
\end{equation}

Dove la corrente entrante $i_{Z_1}$ si ripartisce tra il ramo di filtraggio diretto verso massa ($i_{Z_{24}}$) e il ramo di retroazione positiva ($i_{Z_3}$). Applicando la \textbf{Legge di Ohm generalizzata} alle singole impedenze, esprimiamo le correnti in funzione dei potenziali di nodo:

\begin{itemize}
    \item \textbf{Corrente $i_{Z_1}$:} Definita dalla differenza di potenziale tra l'ingresso e il nodo interno:

    \begin{equation}
        i_{Z_1} = \frac{v_{in} - v_x}{Z_1}
        \label{eq: i_Z1}
    \end{equation}

    \item \textbf{Corrente $i_{Z_2}$:} Definita dalla differenza di potenziale tra il nodo interno e il morsetto non invertente dell'OpAmp:

    \begin{equation}
        i_{Z_2} = \frac{v_x - v^+}{Z_2}
        \label{eq: i_Z2}
    \end{equation}

    \item \textbf{Corrente $i_{Z_3}$:} Corrente che scorre attraverso il ramo di feedback tra il nodo $v_x$ e l'uscita $v_{out}$:

    \begin{equation}
        i_{Z_3} = \frac{v_x - v_{out}}{Z_3}
        \label{eq: i_Z3}
    \end{equation}
\end{itemize}
Sostituendo le espressioni \ref{eq: i_Z1}, \ref{eq: i_Z2} e \ref{eq: i_Z3} nell'equazione di equilibrio \ref{eq: LKT V_x}, otteniamo la relazione fondamentale del sistema:

\begin{equation}
    \frac{v_{in} - v_x}{Z_1} = \frac{v_x-v^+}{Z_2} + \frac{v_x - v_{out}}{Z_3}
    \label{eq: LKT completo}
\end{equation}

Questa equazione presenta tre incognite: $v_x,\, v^+$ e $v_{out}$. Per determinare quindi la funzione di trasferimento, dobbiamo ricorrere al \textbf{principio del corto circuito virtuale:} in un circuito con OpAmp retroazionato negativamente, la tensione sul morsetto invertente tenderà ad essere uguale alla tensione sul morsetto non invertente. 
Per cui, dato che è presente una retroazione negativa, si ricava facilmente che \(v^+ = v^- =v_{out}\).
\begin{equation}
        \frac{v_{in} - v_x}{Z_1} = \frac{v_x-v_{out}}{Z_2} + \frac{v_x - v_{out}}{Z_3}
        \label{eq: LKT con vout}
\end{equation}
Come ultimo passaggio dobbiamo considerare il morsetto non invertente. Applicando la LKC al nodo $v^+$ (ricordando sempre che si considera l'OpAmp ideale, quindi non entrano correnti dai morsetti) notiamo che la corrente che circola sull'impedenza $Z_2$ deve essere la stessa che circola su $Z_4$. Per cui:
\begin{equation}
    i_{Z_2} = i_{Z_4} 
    \label{eq: iz2 iz4}
\end{equation}
Ricordando sempre che $v^+ = v_{out}$ per il principio del c.c. virtuale, queste correnti possono essere trovate applicando la legge di Ohm:
\begin{equation}
    \frac{v_x-v_{out}}{Z_2} = \frac{v_{out}}{Z_4}
    \label{eq: Ohm Z2 Z4}
\end{equation}
Con pochi passaggi algebrici si trova la relazione lineare che lega le tensioni $v_x$ e $v_{out}$:

\begin{equation}
v_x = v_{out}\left(\frac{Z_2}{Z_4}+1\right)
\label{eq: vx vout}  
\end{equation}

Unendo quindi l'equazione \ref{eq: LKT con vout} e l'equazione \ref{eq: vx vout}:
\begin{equation}
    \frac{v_{in} - v_{out}\left(\frac{Z_2}{Z_4}+1\right) }{Z_1} = \frac{v_{out}\left(\frac{Z_2}{Z_4}+1\right) -v_{out}}{Z_2} + \frac{v_{out}\left(\frac{Z_2}{Z_4}+1\right)  - v_{out}}{Z_3}
    \label{eq: LKT vin vout}
\end{equation}

Con pochi passaggi algebrici è possibile definire la funzione di trasferimento di un generico filtro Sallen-Key:

\begin{equation}
    H(j\omega) \triangleq \frac{v_{out}}{v_{in}} = \frac{Z_3Z_4}{Z_1Z_2+Z_3(Z_1+Z_2)+Z_3Z_4}
    \label{eq: Guadagno Sallen Key}
\end{equation}







\section{Filtro Sallen-Key a guadagno non unitario}

\SallenKeyK

In questa configurazione, il guadagno del filtro è determinato dai resistori $R_a$ e $R_b$ che formano un partitore di tensione tra l'uscita e il morsetto non invertente dell'OpAmp.

Per ricavare la funzione di trasferimento si seguono gli stessi passaggi della sezione precedente, con l'unica differenza che ora il principio del corto circuito virtuale ci dice che 
\begin{equation}
    v^+ = v^- = v_y = v_{out}\cdot\frac{R_b}{R_a + R_b}
    \label{eq: v_out partitore}
\end{equation}

Quindi possiamo riprendere l'equazione \ref{eq: LKT con vout} e sostituire \(v^+\) con \(v_{out}\frac{R_b}{R_a + R_b}\):
\begin{equation}
        \frac{v_{in} - v_x}{Z_1} = \frac{v_x-v_{out}\frac{R_b}{R_a + R_b}}{Z_2} + \frac{v_x - v_{out}}{Z_3}
        \label{eq: vx vout partitore}
\end{equation}
Riprendendo l'equazione \ref{eq: Ohm Z2 Z4} e l'equazione \ref{eq: v_out partitore}, possiamo trovare la relazione lineare che lega le tensioni $v_x$ e $v_{out}$:
\begin{equation}
    \frac{v_x-v_{out}\frac{R_b}{R_a + R_b}}{Z_2} = \frac{v_{out}\frac{R_b}{R_a + R_b}}{Z_4} \quad\implies\quad v_x = v_{out}\frac{R_b}{R_a + R_b}\left(\frac{Z_2}{Z_4}+1\right)
    \label{eq: vx con guadagno}
\end{equation}

Sostituendo quindi l'equazione \ref{eq: vx con guadagno} nell'equazione \ref{eq: vx vout partitore} e svolgendo pochi passaggi algebrici, si trova la funzione di trasferimento del filtro Sallen-Key a guadagno non unitario:
\begin{equation}
    H(j\omega) = \frac{v_{out}}{v_{in}} = \frac{Z_3Z_4}{Z_1Z_2+Z_3(Z_1+Z_2)+Z_3Z_4\left(1+\frac{R_b}{R_a}\right)}
    \label{eq: Guadagno Sallen Key non unitario}
\end{equation}

Nel filtro Sallen-Key l'amplificatore operazionale è configurato come
amplificatore non invertente con guadagno $1+\frac{R_b}{R_a}$.
Tale fattore non compare come guadagno moltiplicativo dell'intera
funzione di trasferimento, poiché l'operazionale amplifica il nodo
interno della rete di impedenze e l'uscita viene reintrodotta nella
rete passiva tramite retroazione. 

Di conseguenza il termine
$1+\frac{R_b}{R_a}$ compare nel denominatore della funzione di
trasferimento, contribuendo al polinomio caratteristico del filtro.
Esso quindi non modifica uniformemente il guadagno del sistema, ma
influenza la posizione dei poli e il fattore di qualità $Q$.








